\section{Compare classification branches with a table}\label{utf8-measurement}
The alternative implementation replaces the idle-state multibyte classifier with
a 256-entry lookup table. Each four-byte entry packs a continuation count and two
next-byte bounds, or marks an invalid lead. The data occupy 1024 bytes. ASCII still
takes the direct completion path, and pending continuations still use the retained
state bounds. A lookup therefore affects lead classification rather than every
input byte.

\Needspace{13\baselineskip}
\begin{longtable}{lrrrr}
\toprule Fixture&Branch instr.&Table instr.&Branch reads&Table reads\\\midrule
64 ASCII bytes&991&991&71&71\\
32 two-byte scalars&1567&1791&71&103\\
16 four-byte scalars&1631&1583&71&87\\
Bad \texttt{e0 9f}&64&64&9&10\\
Final \texttt{f0 90}&91&88&9&10\\\bottomrule
\end{longtable}
Each measurement starts with an empty active state, supplies a single final chunk,
and counts completed native instructions through the function return, stopping
before the continuation. Setup, assembly, session construction, main, and exit are
excluded. The first three chunks are respectively 64 copies of \texttt{41}, 32
copies of \texttt{c3 a9}, and 16 copies of \texttt{f0 9f 99 82}. The last two
workloads exercise a detected byte error and a possible prefix truncated at final
input. Both implementations return the same full state and accepted count.

Reads count simulator data-access notices in the observed data region: seven state
word reads, input byte reads, and any metadata-table word reads. They exclude
instruction fetches and external inspection. They are neither a byte total nor a
cache-miss count. The table implementation adds 32 word reads for the two-byte
workload and sixteen for the four-byte workload; it adds none for ASCII.

The table reduces this implementation's instruction count for the four-byte class,
but increases it for the two-byte class. Its address calculation, load, and packed
field extraction cost more than the short two-byte branch path. The four-byte
branch path has additional distinctions for \texttt{f0} and \texttt{f4}, where the
table instead returns the prepared bounds. The examples identify a workload-dependent
tradeoff rather than a universal table-lookup speedup.

A later optimization might scan a bounded ASCII span or use a smaller classification
table. Preserve pending-state handling, stop at the first impossible byte, and keep
full access widths inside each chunk. Native instruction counts alone do not
establish host elapsed time or physical processor cycles. Measure the relevant
quantity separately after the complete behavioral checks still pass.
